\section{Integration and positivity for finite residue words}
\label{ji-section}

The first-coordinate integration can be performed for every finite
matrix of polynomial differential operators at once.  The essential
observation is that the scalar Jacquet marginals already form a
matrix of Schwartz functions.  Thus the derivative terms in the
two-component identity have zero integral without a separate decay
argument for each subset of a word.  A second argument extends any
positivity proved for compact coefficients to a prescribed class with
finitely many decay and differentiability bounds.  An ordered matrix
identity handles both residue operations together, giving actual local
positivity for every fixed finite word under the terminal hypotheses.

Throughout this section the place data, Haar measures, compact
representations and positive Gram support are those of
Section~\ref{jf-section}.  The standard flag and its normalizations
are those of Corollary~\ref{cm-chain-restriction}.  All ranks, words
and spectral parameters are fixed unless a uniform range is stated.

\subsection{A matrix of Schwartz marginals}

Write
\[
 \mathscr C_\varepsilon(G_r,E)
   =\{F\in C_c^\infty(G_r,E):F=0\text{ on }
                                      \mathcal G_r<\varepsilon\}.
\]
This is a subspace of \(\mathscr S_\varepsilon\).
Let \(B\) be a finite auxiliary index set.  For vectors indexed by
\(B\), the Jacquet transform acts on each component, and all inner
products are summed over \(B\).  The inner product is linear in its
first argument.  For \(r\geq1\), put
\(\mathcal H=L^2(d\eta_{r-1},E)\).

Let \(\mathcal Q_r(Y)\) denote the full contracted matrix on
\(G_r\), with the real spectral normalization of
Corollary~\ref{cm-group-preservation}.  An admissible polynomial
matrix in this section means a finite matrix whose entries are
finite sums of ordered products of entries of
\(\mathcal Q_r(Y)\), evaluated at fixed real parameters, with
constant scalar or auxiliary-matrix coefficients.  The products
need not commute.  Constant denominators in the parameters are
allowed when nonzero; inverse coordinate, Gram or differential
operators are not included in this definition.

\begin{theorem}\label{ji-schwartz-matrix}
Let \(r\geq1\), \(F,H\in\mathscr C_\varepsilon(G_r,E)^B\), and let
\(P\) be an admissible polynomial matrix acting on these vectors.
Set
\[
 U(t)=\Phi_{PF}(t,\cdot),\qquad
 V(t)=\Phi_H(t,\cdot)\qquad(t\ne0).
\]
Then \(U,V:\mathbb C^\times\to\mathcal H^B\) are smooth.
Every entry of the finite matrix
\begin{equation}\label{ji-kernel}
 K_{ab}(t)=|t|^{-2}\langle U_a(t),V_b(t)\rangle_{\mathcal H}
 \qquad(a,b\in B)
\end{equation}
has a unique Schwartz extension to \(\mathbb C\).
Consequently, for every constant matrix \(A\) on the auxiliary
space,
\begin{equation}\label{ji-trace-density}
 D(U,AV)(t)=|t|^{-2}\langle U(t),AV(t)\rangle
                         =\operatorname{tr}(K(t)A^*)
\end{equation}
has a Schwartz extension.  Every positive-order constant-coefficient
derivative of this extension has integral zero over \(\mathbb C\).
\end{theorem}

\begin{proof}
Corollary~\ref{cm-group-preservation} shows that every component of
\(PF\) is smooth, compactly supported, and zero on the same open
set \(\mathcal G_r<\varepsilon\).  Corollary~\ref{jf-smooth-slices}
therefore gives the Hilbert-space smoothness.  Apply
Theorem~\ref{jf-marginal-schwartz} to each pair
\(((PF)_a,H_b)\).  It gives precisely the pointwise formula
\eqref{ji-kernel} on \(\mathbb C^\times\), as well as its Schwartz
extension across zero.  There are only finitely many pairs.

Linearity in the first argument gives
\(D(U,AV)=\sum_{a,b}K_{ab}\overline{A_{ab}}\), which is
\eqref{ji-trace-density}.  A derivative of a Schwartz function is
integrable.  Integrating one real coordinate derivative first gives
zero, since the function tends to zero at both ends; Fubini is
valid by absolute integrability.  Repeating this observation proves
the assertion for all positive-order derivatives, including the
Wirtinger derivatives used below.
\end{proof}

\begin{corollary}\label{ji-integrated-word}
Take \(B=\{0,1\}^S\), and let
\(w=((a_j,\sigma_j,Y_j))_{j=1}^h\) have pairwise distinct
\(a_j\in S\), signs \(\sigma_j\in\{+,-\}\), and real
parameters \(Y_j\).  Use the ordered operators \(I_w,V_w\)
of Section~\ref{rf-section}, with the rightmost factor acting first.
For \(U,V\) in Theorem~\ref{ji-schwartz-matrix}, the pointwise
density \(D(I_wU,V_wV)\) has a Schwartz extension and
\begin{equation}\label{ji-integrated-identity}
 \int_{\mathbb C}D(I_wU,V_wV)(t)\,dt
                         =\int_{\mathbb C}D(U,V)(t)\,dt.
\end{equation}
\end{corollary}

\begin{proof}
Proposition~\ref{rf-subset-expansion} gives, on
\(\mathbb C^\times\),
\[
 D(I_wU,V_wV)=\sum_{J\subseteq\{1,\ldots,h\}}
                    \epsilon_J\delta_JD(U,E_JV).
\]
Every \(E_J\) is a constant matrix.  Theorem~\ref{ji-schwartz-matrix}
makes every term on the right a Schwartz function.  Terms with
\(J\ne\varnothing\) have zero integral by that theorem, and the
empty term is \(D(U,V)\).  This proves the asserted extension and
the equality.  The conclusion concerns the scalar paired density;
it does not assert that either transformed vector separately belongs
to \(L^2(|t|^{-2}dt;\mathcal H^B)\).
\end{proof}

The empty word gives the identity.  Distinct auxiliary indices
remain essential: no decay statement repairs an algebraically
incorrect repeated-index expansion.  The constants implicit in
Schwartz seminorms may depend on the whole fixed word and on the
coefficients.  No estimate uniform over unbounded ranks or word
lengths is asserted.

\subsection{Positivity reduced to an algebraic congruence}

For an admissible polynomial matrix \(P_j\) on \(G_j\), let
\(P_j^{\rm rad}\) be its radial restriction.  Entrywise
intertwining in Corollary~\ref{cm-chain-restriction} and composition
give
\begin{equation}\label{ji-polynomial-intertwining}
          P_j^{\rm rad}\Phi_F=\Phi_{P_jF}.
\end{equation}

\begin{corollary}\label{ji-word-pairing}
Fix the place data and Haar measures of Section~\ref{jf-section}.
There is a constant \(C_r>0\), independent of the auxiliary dimension
and of the admissible polynomial matrix \(P\), such that for all
\(F,H\in\mathscr S_\varepsilon(G_r,E)^B\),
\[
 \sum_{a\in B}\beta((PF)_a,H_a)
   =C_r^{-1}\langle P^{\mathrm{rad}}\Phi_F,\Phi_H\rangle_{L^2(d\eta_r)}.
\]
In particular the same constant works for every finite residue word
whose scalar denominators are nonzero.
\end{corollary}

\begin{proof}
The coefficient class is preserved by \(P\), by
Corollary~\ref{cm-group-preservation}. Apply
\eqref{jf-radial-pairing} componentwise to \(PF,H\), and use
\eqref{ji-polynomial-intertwining}. Finite component sums preserve the
same constant. The intertwining identity holds for any ordered polynomial
because it holds for each contracted entry and passes to sums and
compositions. Constant auxiliary matrices, including insertion and
transposition at a fresh auxiliary factor, commute with the componentwise
Jacquet transformation. Thus it applies to both signed residue steps.
\end{proof}

Define its quadratic form on compact coefficient vectors by
\[
 \mathfrak q_j(F)=\int
       \langle P_j^{\rm rad}\Phi_F,\Phi_F\rangle\,d\eta_j.
\]
This integral is absolutely convergent, by
\eqref{ji-polynomial-intertwining} and Cauchy--Schwarz.

\begin{theorem}\label{ji-rank-descent}
Fix \(r\), a finite auxiliary set \(S\), and admissible polynomial
matrices \(P_j\), \(0\leq j\leq r\), on the corresponding
standard flag.  Suppose that for each \(j\geq1\) there is a word
\(w_j\) with distinct auxiliary indices such that, as differential
operators on the open radial domain,
\begin{equation}\label{ji-congruence}
 P_j^{\rm rad}=V_{w_j}^*P_{j-1}^{\rm rad}V_{w_j}.
\end{equation}
Here \(V_{w_j}\) acts in the outermost coordinate; the star is the
formal adjoint for its measure \(|t|^{-2}dt\).
If \(\mathfrak q_0(F)\geq0\) for every compact smooth terminal
coefficient vector with a positive Gram cutoff, then
\(\mathfrak q_r(F)\geq0\) for every
\(F\in\mathscr C_\varepsilon(G_r,E)^{\{0,1\}^S}\).

More precisely, writing \(F_t\) for the compact family of
Proposition~\ref{jf-compact-family} and applying \(V_{w_j}\)
to its parameter \(t\), one has
\begin{equation}\label{ji-positive-descent}
 \mathfrak q_j(F)=
   \int_{\mathbb C}|t|^{-2}
                    \mathfrak q_{j-1}(V_{w_j}F_t)\,dt.
\end{equation}
The scalar integrand has a Schwartz extension after the factor
\(|t|^{-2}\) is included.
\end{theorem}

\begin{proof}
At each nonzero \(t\), every component of \(V_{w_j}F_t\) is a
finite linear combination of parameter derivatives of \(F_t\),
multiplied by smooth scalar functions of \(t\).  On a compact
neighborhood of this parameter, Proposition~\ref{jf-compact-family}
gives common compact support and a common positive Gram cutoff.
Thus \(V_{w_j}F_t\) is an admissible compact coefficient vector
on the smaller group.  Its Jacquet transform equals
\(V_{w_j}\Phi_F(t,\cdot)\), by the same proposition and
Corollary~\ref{jf-smooth-slices}.

Put \(U=\Phi_{P_jF}\) and \(V=\Phi_F\).  The formal inverse
identity of Lemma~\ref{rf-inverse-adjoint} and
\eqref{ji-congruence} imply, on this smooth domain,
\[
 I_{w_j}U=P_{j-1}^{\rm rad}V_{w_j}\Phi_F.
\]
This is an identity of differential expressions, so no domain of a
closed unbounded operator is being assumed.  Apply
Corollary~\ref{ji-integrated-word}.  Its right side is
\(\mathfrak q_j(F)\), and its left side is the right side of
\eqref{ji-positive-descent}.  Each inner smaller-rank pairing
exists: both vectors are Jacquet transforms of compact coefficient
vectors, one obtained by applying \(P_{j-1}\).
The scalar extension supplied by that corollary proves absolute
integrability in \(t\); no exchange of a merely conditionally
convergent multiple integral has been used.

If \(\mathfrak q_{j-1}\) is nonnegative on all such compact
coefficients, the integrand is nonnegative for every \(t\ne0\).
Equation~\eqref{ji-positive-descent} proves the same property for
\(\mathfrak q_j\).  Induction from \(j=0\) proves the theorem.
\end{proof}

\begin{corollary}\label{ji-one-insertion}
For a real \(Y\) and \(\sigma\in\{1,-1\}\), consider either
\(P_j=\sigma\mathcal Q_j(Y)\) for all \(j\), or
\(P_j=\sigma\mathcal Q_j(Y)^t\) for all \(j\), where the
transpose is on the two auxiliary indices.  If the corresponding
terminal compact form is nonnegative, the actual Jacquet form is
nonnegative at every finite rank.

In the notation of the initial frame of
Section~\ref{re-section}, sufficient terminal conditions are
\[
\begin{array}{ll}
\text{plus, }Y=a:
 &a,a+1\in I,\quad q(a)>0,\quad
           \sigma p_\rho(a)>0,\quad\sigma p_\rho(a+1)>0;\\[2pt]
\text{minus, }Y=b:
 &b-1,b,b+1\in I,\quad q(b-1)>0,\quad
           \sigma p_\rho(b-1)>0,\quad\sigma p_\rho(b)>0.
\end{array}
\]
The inequalities are required for every terminal finite type
\(\rho\).  They make the two weights in
Corollary~\ref{re-density} positive.  In the minus case the
conditions use the shifted parameter \(y=b-1\), whereas the
radial operator is \(V_-(b)\).  Positivity of the contracted
matrix at an arbitrary real parameter is not asserted.
\end{corollary}

\begin{proof}
Write \(p=t\partial_t\), \(\bar p=\bar t\partial_{\bar t}\),
and let \(A(Y)\) be the contracted matrix in the remaining
coordinates.  The one-pair restriction is
\[
 \mathcal Q_j^{\rm rad}(Y)=L(Y)A(Y)T(Y),\quad
 T(Y)=\begin{pmatrix}Y+p&-t\\t^{-1}&0\end{pmatrix},\quad
 L(Y)=\begin{pmatrix}Y-\bar p&\bar t^{-1}\\-\bar t&0\end{pmatrix}.
\]
For \(|t|^{-2}dt\), one has \(p^*=-\bar p\).  Hence
\(L(Y)=T(Y)^*\), and \(T(Y)=V_+(Y)\).  This proves
\eqref{ji-congruence} for the plus case.

For the transpose, let
\[
 W(Y)=V_-(Y)=
       \begin{pmatrix}Y-\bar p&-\bar t\\\bar t^{-1}&0\end{pmatrix}.
\]
The entries of \(A(Y)\) commute with all outer coordinate
operators.  Every entry of \(T(Y)\) also commutes with every
entry of \(L(Y)\): the only Euler derivative in \(T\) is
\(p\), acting on holomorphic coordinates, and the only Euler
derivative in \(L\) is \(\bar p\), acting on antiholomorphic
coordinates.  Matrix multiplication can therefore be transposed
in the ordinary reversed order in this particular product.  Since
\(L(Y)^t=W(Y)\) and \(T(Y)^t=W(Y)^*\), it follows that
\[
       (L(Y)A(Y)T(Y))^t=W(Y)^*A(Y)^tW(Y).
\]
No commutation between two entries of \(A\) is required.
Multiplication by the real sign \(\sigma\) preserves both
congruences.  Theorem~\ref{ji-rank-descent} now applies.  At rank
zero the compact-support statement at the end of
Section~\ref{jf-section} makes the Jacquet vector an admissible
compact terminal vector for Corollary~\ref{re-density}.
\end{proof}

\subsection{Arbitrary residue words}
\label{mx-section}

Label the fresh two-dimensional auxiliary factors chronologically
by \(1,\ldots,h\).  On factors \(i,j\), put
\[
 \Pi_{ij}=\sum_{a,b=0}^1 E_{ab}^{(i)}E_{ba}^{(j)},\qquad
 \Omega_{ij}=\sum_{a,b=0}^1 E_{ab}^{(i)}E_{ab}^{(j)},
\]
and define
\[
 R_{\Pi,ij}(z)=\frac{z\,1-\Pi_{ij}}{z-1},\qquad
 R_{\Omega,ij}(z)=\frac{z\,1-\Omega_{ij}}{z-1}.
\]
These are the permutation and contraction matrices of
\eqref{re-fresh-matrices}.  Transposition in either one of their
two factors exchanges \(\Pi\) and \(\Omega\).

Choose signs \(\epsilon_i\in\{1,-1\}\) and real parameters
\(\lambda_i\).  A plus step means \(\epsilon_i=1\) and
\(a_i=\lambda_i\); a minus step means \(\epsilon_i=-1\)
and \(b_i=\lambda_i\).  For a matrix \(A(Y)\) in the remaining
variables, let \(A_i^{[1]}(Y)=A_i(Y)\) and
\(A_i^{[-1]}(Y)=A_i(Y)^t\), with the indices inserted at factor
\(i\).  Define the constant auxiliary matrices
\begin{align}
 C_i&=\prod_{\substack{j<i\\\epsilon_j\ne\epsilon_i}}^{\text{decreasing }j}
       R_{\Omega,ij}\bigl(\epsilon_i(\lambda_j-\lambda_i)\bigr),
                                                        \label{mx-C}\\
 D_i&=\prod_{\substack{j<i\\\epsilon_j=\epsilon_i}}^{\text{increasing }j}
       R_{\Pi,ij}\bigl(\epsilon_i(\lambda_i-\lambda_j)\bigr).
                                                        \label{mx-D}
\end{align}
Products are written in the indicated order, and an empty product
is the identity.  Assume every displayed denominator is nonzero.

\begin{lemma}\label{mx-recursion}
For the actual mixed recursion of
\eqref{re-plus-module} and \eqref{re-minus-module}, starting
from \(\mathcal R_0(A)=1\), its residue matrix is exactly
\begin{equation}\label{mx-residue-recursion}
 \mathcal R_i(A)=\mathcal R_{i-1}(A)
                    C_i A_i^{[\epsilon_i]}(\lambda_i)D_i.
\end{equation}
In particular it is a polynomial of degree at most \(i\) in
the entries of the original \(A\), with scalar spectral
denominators and constant auxiliary coefficients.
\end{lemma}

\begin{proof}
After the first \(i-1\) operations, the main matrix is
\[
 \left(\prod_{\substack{j<i\\\epsilon_j=-1}}^{\text{decreasing }j}
                         R_Q(\lambda_j-Y)\right)
 A(Y)
 \left(\prod_{\substack{j<i\\\epsilon_j=1}}^{\text{increasing }j}
                         R_P(Y-\lambda_j)\right).
\]
This follows directly because a minus operation multiplies on the
left and a plus operation on the right.  A new plus insertion
evaluates this expression at \(Y=\lambda_i\) and places the
main index at factor \(i\), giving \(C_i A_i D_i\).

A new minus insertion evaluates and transposes the main index.
This reverses the written product, changes \(R_P\) to
\(R_\Omega\) on the left, and changes \(R_Q\) to
\(R_\Pi\) on the right.  This particular partial transpose is
legitimate: the coefficient entries from distinct old auxiliary
factors commute with one another and with the entries of the
original \(A\).  It gives precisely
\(C_i A_i^tD_i\), with the orders in \eqref{mx-C} and
\eqref{mx-D}.  Multiplication by the previous residue proves
\eqref{mx-residue-recursion}.  Each step introduces one entry of
the original \(A\), proving the degree bound.
\end{proof}

\subsection{The mixed congruence}

Use the outer-coordinate operators
\[
 T(Y)=\begin{pmatrix}Y+p&-t\\t^{-1}&0\end{pmatrix},\qquad
 W(Y)=\begin{pmatrix}Y-\bar p&-\bar t\\\bar t^{-1}&0\end{pmatrix},
 \qquad p=t\partial_t,\quad\bar p=\bar t\partial_{\bar t}.
\]
Formal adjoints are for \(|t|^{-2}dt\), so
\(W(Y)^*=T(Y)^t\) and \(T(Y)^*=W(Y)^t\) for real \(Y\).
Write
\[
 U_i=
 \begin{cases}
 T_i(\lambda_i),&\epsilon_i=1,\\
 W_i(\lambda_i),&\epsilon_i=-1,
 \end{cases}
 \qquad V_i=U_i\cdots U_1.
\]


\begin{lemma}\label{ji-transfer-relations}
For distinct auxiliary factors and real parameters where the rational
matrices are defined,
\begin{align}
 T_j(u)T_i(v)R_{\Pi,ij}(v-u)
       &=R_{\Pi,ij}(v-u)T_i(v)T_j(u),\label{ji-transfer-rtt}\\
 W_j(u)W_i(v)R_{\Pi,ij}(u-v)
       &=R_{\Pi,ij}(u-v)W_i(v)W_j(u).\label{ji-opposite-rtt}
\end{align}
\end{lemma}

\begin{proof}
Using matrix indices \(0,1\), write \(T(Y)=Y E_{00}+S\).
Substitution in the component
relations \eqref{re-rtt}, as a polynomial identity in \(u,v\),
reduces them to
\[
 [S_{ab},S_{cd}]=\delta_{c0}\delta_{b0}S_{ad}
                              -\delta_{a0}\delta_{d0}S_{cb}.
\]
Their nonzero cases are precisely
\([p,-t]=-t\) and \([p,t^{-1}]=-t^{-1}\), together with
antisymmetry; the remaining cases use \([-t,t^{-1}]=0\)
and \(S_{11}=0\).  Multiplication by the permutation matrix
writes these relations as \eqref{ji-transfer-rtt}.  Multiplying
by its rational scalar normalization does not change the identity.

For \(W\), the corresponding commutators are
\([-\bar p,-\bar t]=\bar t\) and
\([-\bar p,\bar t^{-1}]=\bar t^{-1}\), with commuting
off-diagonal entries.  They negate the preceding component
relations, replacing \(v-u\) by \(u-v\).  This proves
\eqref{ji-opposite-rtt}.  The argument uses polynomial identities
before division, so no extra exclusion of \(u=v\) is needed.
\end{proof}

\begin{lemma}\label{mx-transfer-moves}
With \(C_i,D_i\) as above,
\begin{equation}\label{mx-two-moves}
          V_{i-1}C_iU_i^*=U_i^*C_iV_{i-1},
       \qquad V_{i-1}U_iD_i=D_iU_iV_{i-1}.
\end{equation}
\end{lemma}

\begin{proof}
Use the two quadratic identities in
Lemma~\ref{ji-transfer-relations}.
Partially transpose factor \(i\) in the first identity.  Keep
\(T_j(u)\) on its original side: it has no matrix index in
factor \(i\).  Only the products \(T_iR_\Pi\) and
\(R_\Pi T_i\) reverse their order under this operation.
Their coefficient algebras commute, since the entries of \(T_i\)
act on the outer coordinate and the coefficients of \(R_\Pi\)
act on factor \(j\).  Thus the resulting equality is
\[
 T_j(u)R_{\Omega,ij}(v-u)W_i(v)^*
                  =W_i(v)^*R_{\Omega,ij}(v-u)T_j(u).
\]
This does not commute the noncommuting entries of \(T_i\)
and \(T_j\).  The identical partial transpose in the second
quadratic identity gives
\[
 W_j(u)R_{\Omega,ij}(u-v)T_i(v)^*
                  =T_i(v)^*R_{\Omega,ij}(u-v)W_j(u).
\]
These are exactly the required relations for old and fresh factors
of opposite signs, with the arguments in \eqref{mx-C}.

To obtain the first identity in \eqref{mx-two-moves}, process the
old factors in increasing order of their index.  An opposite-sign
factor passes the fresh adjoint by the corresponding relation just
proved.  A same-sign factor commutes with that adjoint, since one
uses only holomorphic outer operators and the other only
antiholomorphic ones.  Constant matrices involving other old
indices commute with the factor currently being moved.  The
decreasing order in \(C_i\) is therefore preserved throughout.

For the second identity, a same-sign old factor uses the original
quadratic relation, and an opposite-sign old factor commutes with
the fresh \(U_i\).  Again the unused constant factors act on
disjoint old indices.  Processing old indices in increasing order
leaves the increasing order of \(D_i\) unchanged.  This proves
both displayed identities without exchanging two old transfer
factors.
\end{proof}

\begin{theorem}\label{ji-fused-congruence}
Let the entries of \(A(Y)\) commute with all outer-coordinate
operators, and put \(Q(Y)=T(Y)^*A(Y)T(Y)\).  For every fixed
finite word of plus and minus insertions with fresh auxiliary
factors and nonzero scalar denominators,
\begin{equation}\label{mx-full-congruence}
              \mathcal R_h(Q)=V_h^*\mathcal R_h(A)V_h.
\end{equation}
The entries of \(A\) need not commute with one another.
\end{theorem}

\begin{proof}
Corollary~\ref{ji-one-insertion} gives
\[
       Q_i^{[\epsilon_i]}(\lambda_i)
            =U_i^*A_i^{[\epsilon_i]}(\lambda_i)U_i.
\]
Assume \eqref{mx-full-congruence} at length \(i-1\), and
abbreviate \(A_i^{[\epsilon_i]}(\lambda_i)\) to \(A_i'\).
Lemma~\ref{mx-recursion} and the two ordered moves yield
\[
\begin{split}
 \mathcal R_i(Q)
 &=V_{i-1}^*\mathcal R_{i-1}(A)V_{i-1}
                              C_iU_i^*A_i'U_iD_i\\
 &=V_{i-1}^*U_i^*\mathcal R_{i-1}(A)C_i
                                      A_i'V_{i-1}U_iD_i\\
 &=V_i^*\mathcal R_{i-1}(A)C_iA_i'D_iV_i\\
 &=V_i^*\mathcal R_i(A)V_i.
\end{split}
\]
In the second line, \(\mathcal R_{i-1}(A)\) commutes with
\(U_i^*\), since factor \(i\) is fresh and the remaining
variables are independent of the outer coordinate.  Also
\(A_i'\) commutes with \(V_{i-1}\), whose auxiliary factors
are all old.  No commutation of the previous residue with
\(C_i\), or between two entries of \(A\), is used.
The identity at length zero is the identity matrix, completing
the induction.
\end{proof}


\begin{corollary}\label{ji-single-sign-congruence}
For a word with only plus insertions, or only minus insertions,
respectively, the general congruence specializes to
\begin{align}
 \mathcal R_h^+(Q)
   &=\mathcal V_+^*\mathcal R_h^+(A)\mathcal V_+,
 &\mathcal V_+&=T_h(a_h)\cdots T_1(a_1),
                                      \label{ji-plus-congruence}\\
 \mathcal R_h^-(Q)
   &=\mathcal V_-^*\mathcal R_h^-(A)\mathcal V_-,
 &\mathcal V_-&=W_h(b_h)\cdots W_1(b_1).
                                      \label{ji-minus-congruence}
\end{align}
The rightmost factor acts first.
\end{corollary}

\begin{proof}
Take all signs equal in Theorem~\ref{ji-fused-congruence}.
Then every \(C_i\) is the identity and the same ordered formula
applies.
\end{proof}

\subsection{One passage from compact support to finite decay}

Let \(d=k+2r+1\) be the dimension of the defining Hermitian space,
and let \(M(g)=\|g\|_{\mathrm F}\), as in
Section~\ref{jm-section}.  In particular \(M\geq1\).
For an auxiliary vector, use the product Hilbert norm in the
seminorm \(\|\cdot\|_{A,m}\) of Section~\ref{jp-section}.

\begin{lemma}\label{ji-degree-order}
An admissible polynomial matrix \(P\) of degree at most \(h\)
has differential order at most \(dh\).  In an invariant
derivative expansion, its coefficients, and all their fixed mixed
left and right derivatives, are bounded by constants times
\(M^{2h}\).  Consequently, for every integer \(m\geq0\),
\begin{equation}\label{ji-operator-bound}
       \|PF\|_{A,m}\leq C_{P,m}\|F\|_{A+2h,dh+m}.
\end{equation}
The constant is finite for fixed parameters; it can depend on their
denominators and on the size of the auxiliary space.
\end{lemma}

\begin{proof}
The full Capelli determinant is a polynomial of order at most
\(d\) in invariant derivatives, with constant coefficients.
The coefficient \(Z(g)\) and its conjugate are restrictions of
polynomials of degree one in matrix entries and their conjugates.
Invariant differentiation preserves this polynomial degree.
Expanding a commutator with multiplication by \(Z\) lowers
differential order by at least one, and gives coefficients of
degree at most one.  A second commutator with \(\bar Z\)
lowers order once more and gives degree at most two.
The four entries in \eqref{cm-four-commutators} therefore all
have order at most \(d\) and coefficient degree at most two.

In a product of \(h\) entries, the Leibniz rule differentiates
coefficients without increasing their degree.  Its order is at most
\(dh\), and its coefficient degree at most \(2h\).
The same degree bound holds after any fixed invariant derivatives
of the coefficients.  A polynomial of degree at most \(2h\)
is bounded by a constant times \(M^{2h}\), since \(M\geq1\).
Mixed left and right derivatives of the output give only finitely
many terms of this kind and derivatives of \(F\) of total order
at most \(dh+m\); left and right invariant derivatives commute.
This proves \eqref{ji-operator-bound}, also for finite sums and
constant auxiliary matrices.
\end{proof}

\begin{lemma}\label{ji-compact-cutoff}
Fix \(\delta>0\), \(A\geq0\), and \(m\geq0\).  There are
smooth compactly supported scalar functions \(\chi_R\), \(R\geq1\),
such that \(\chi_RF\in\mathscr C_\varepsilon\) whenever
\(F\) is smooth and vanishes on \(\mathcal G_r<\varepsilon\),
and
\begin{equation}\label{ji-cutoff-bound}
 \|(1-\chi_R)F\|_{A,m}
                   \leq C_{m,\delta}R^{-\delta}
                                      \|F\|_{A+\delta,m}.
\end{equation}
The support of \(\chi_R\) is contained in \(M\leq2R\).
\end{lemma}

\begin{proof}
Choose a smooth real function \(\chi\) on \([0,\infty)\)
equal to one on \([0,1]\), zero on \([2,\infty)\), and
between zero and one elsewhere, and put
\(\chi_R(g)=\chi(M(g)/R)\).
The matrix group is closed in matrix space, so this support is
compact.  The function \(M\) is smooth on the group because
it never vanishes.  Every fixed mixed invariant derivative of
\(M^2\) is bounded by a constant times \(M^2\): matrix
multiplication makes the differentiated entries linear in the
original entries.  The chain rule, applied successively to the
positive square root, gives the corresponding bound \(CM\)
for every fixed derivative of \(M\).

All positive-order derivatives of \(\chi_R\) are supported
on \(R\leq M\leq2R\) and are bounded independently of
\(R\).  The zero-order function \(1-\chi_R\) vanishes on
\(M<R\) and is bounded by one.  Apply the Leibniz rule to
\((1-\chi_R)F\).  On the support of each resulting term,
\(M^{-\delta}\leq R^{-\delta}\).  This proves
\eqref{ji-cutoff-bound}.  Multiplication does not enlarge the
support of \(F\), so the same Gram cutoff is preserved.
\end{proof}

\begin{theorem}\label{ji-finite-decay-positivity}
Let \(A\geq B(k,r)\), \(\delta>0\), and let \(P\) have
polynomial degree at most \(h\).  For smooth vectors satisfying
\[
             \|F\|_{A+2h,dh+5}<\infty,
             \qquad\|H\|_{A,5}<\infty,
\]
the sesquilinear form
\[
       \mathcal B_P(F,H)=\sum_{a\in B}\beta((PF)_a,H_a)
\]
is well defined and obeys
\begin{equation}\label{ji-form-bound}
 |\mathcal B_P(F,H)|
       \leq C_P\|F\|_{A+2h,dh+4}\|H\|_{A,4}.
\end{equation}
Suppose this form is nonnegative on the diagonal for every
\(F\in\mathscr C_\varepsilon(G_r,E)^B\).
Then it is nonnegative for every smooth \(F\) with the same
positive Gram cutoff and
\begin{equation}\label{ji-finite-class}
                    \|F\|_{A+2h+\delta,dh+5}<\infty.
\end{equation}
Indeed its compact cutoffs satisfy
\begin{equation}\label{ji-form-error}
 |\mathcal B_P(F,F)-\mathcal B_P(\chi_RF,\chi_RF)|
       \leq C_{P,\delta}R^{-\delta}
                             \|F\|_{A+2h+\delta,dh+5}^{\,2}.
\end{equation}
\end{theorem}

\begin{proof}
Lemma~\ref{ji-degree-order} puts \(PF\) in
\(\mathcal C_A^5\), and \(H\) is already in this space.
Theorem~\ref{jp-pairing-equality} therefore defines the actual
Jacquet pairing and bounds it by the fourth-derivative seminorms.
Together with \eqref{ji-operator-bound}, this proves
\eqref{ji-form-bound}; the finite component sum is absorbed in
the constant.

Write \(F_R=\chi_RF\).  Sesquilinearity gives
\[
 \mathcal B_P(F,F)-\mathcal B_P(F_R,F_R)
       =\mathcal B_P(F-F_R,F)+\mathcal B_P(F_R,F-F_R).
\]
Apply \eqref{ji-form-bound} and
Lemma~\ref{ji-compact-cutoff}, first with decay index \(A+2h\)
and then with index \(A\).  The same lemma bounds the norms of
\(F_R\) by fixed multiples of those of \(F\), independently
of \(R\).  Since \(M\geq1\), the norm in
\eqref{ji-finite-class} dominates all the required norms.  This
gives \eqref{ji-form-error}.  Each \(F_R\) belongs to
\(\mathscr C_\varepsilon\), so its diagonal value is real and
nonnegative.  Its limit has the same property.
\end{proof}

For fixed \(k,r,h\), the decay and derivative indices in
\eqref{ji-finite-class} apply to every polynomial residue of
length at most \(h\), at every set of parameters where that
residue is defined.  Uniform operator-norm constants over these
parameters are unnecessary for this positivity conclusion: for
each individual form one takes its compact-support limit.  Such
constants would be necessary for a quantitative estimate uniform
over a varying family.  The theorem supplies neither that estimate
nor uniform control as the ranks and word lengths tend to infinity.

\subsection{Actual local positivity for every finite word}

\begin{corollary}\label{ji-fused-positivity}
Fix \(k\geq1\), \(r\geq0\), the standard flag and the
actual Jacquet place data of Section~\ref{jf-section}.  Fix a
finite word of \(h\) plus and minus insertions, with distinct
auxiliary factors and a single initial sign \(\sigma\).
Assume the initial-domain, shifted-sign and separation hypotheses
of Lemma~\ref{re-shifted-domain}.  Then the residue form of the
explicit recursion in Section~\ref{re-section} is nonnegative on
actual Jacquet transforms of compact smooth coefficient vectors
with positive Gram cutoff.

The same positivity holds for every smooth coefficient vector
with that cutoff and
\[
       \|F\|_{A+2h+\delta,(k+2r+1)h+5}<\infty,
           \qquad A\geq B(k,r),\quad\delta>0.
\]
The assertion is for each fixed finite rank and word.  No uniform
quantitative energy bound over growing ranks or words is claimed.
\end{corollary}

\begin{proof}
At each rank apply Theorem~\ref{ji-fused-congruence} to the
one-pair restriction of \(\sigma\mathcal Q_j\).  The transfer
word \(V_h\) lists the chronological insertions in reverse
order under the product convention of Section~\ref{rf-section};
its auxiliary indices are pairwise distinct.  The terminal residue
has the positive compact form of Corollary~\ref{re-density}.
Theorem~\ref{ji-rank-descent} therefore proves actual positivity
for compact group coefficients, including every mix of signs.

Each smaller-rank pairing in that descent is absolutely defined
for a compact coefficient vector at each nonzero outer parameter.
Its scalar outer density is Schwartz by
Corollary~\ref{ji-integrated-word}.  This justifies the iterated
integration without requiring separate global square integrability
of the transferred vectors or absolute integrability of every
term of a fully expanded multiple integral.

Lemma~\ref{mx-recursion} gives polynomial degree at most \(h\)
in the full contracted entries.  The finite-decay conclusion now
follows directly from Theorem~\ref{ji-finite-decay-positivity},
using the positive normalization constant in
\eqref{jf-radial-pairing}.
\end{proof}

\begin{corollary}\label{ji-single-sign-positivity}
Under the preceding hypotheses, the same conclusion holds in
particular for every fixed finite word containing only plus
insertions or only minus insertions.
\end{corollary}

\begin{proof}
These are the two constant-sign cases of
Corollary~\ref{ji-fused-positivity}.
\end{proof}

\begin{corollary}\label{ji-unshifted-positivity}
Fix the place data, standard flag and initial spectral parameters of
Corollary~\ref{ji-fused-positivity}. Take a finite word on distinct
auxiliary factors with real insertion parameters \(\lambda_j\).
Assume the interlacing and common-cut conditions
\eqref{rp-interlacing} and \eqref{rp-cut-margin}, the latter with
\(Y=\lambda_j\), for one fixed \(1\le l\le k\), and use initial
sign \(\sigma=(-1)^{k-l}\). If
\[
                     |\lambda_i-\lambda_j|>2\qquad(i\ne j),
\]
then the compact and finite-decay positivity conclusions of
Corollary~\ref{ji-fused-positivity} hold. The separation in this
statement is imposed on the original insertion parameters for both signs.
\end{corollary}

\begin{proof}
Lemma~\ref{rp-parameter-signs} gives the required signs of the initial
spectral polynomials at \(\lambda_j-1,\lambda_j,\lambda_j+1\),
and positivity of \(q(\lambda_j)\) and \(q(\lambda_j-1)\).
The original-parameter admissibility argument
\eqref{re-word-domain} makes every recursive step well defined.
Lemma~\ref{rp-recursive-positive} and Corollary~\ref{re-density}
therefore give terminal positivity for every mixture of signs.

Theorem~\ref{ji-fused-congruence} and Theorem~\ref{ji-rank-descent}
need distinct auxiliary factors and nonzero scalar denominators, and
impose no further separation on shifted parameters. Applying them gives
compact-input positivity. Theorem~\ref{ji-finite-decay-positivity}
gives the asserted extension with the same finite decay indices.
\end{proof}

This establishes the local positivity for the explicitly defined
residue recursion and the actual compact or finite-decay Jacquet
coefficients just stated. The integral comparison uses one positive
normalization for every finite word. Global orbital-energy estimates and
the arithmetic nonvanishing and divisibility requirements do not
follow from this local positivity theorem.
